\section{El intercambio de privacidad: ceder context a cambio de conexiones}

El problema de producto de Elys es lograr que los usuarios estén dispuestos a ceder su context. El conductor señala que los usuarios adoptan de forma natural una actitud defensiva: no quieren que una organización comercial sepa demasiado de ellos, no quieren que su doble digital diga cosas indebidas y tampoco tienen claro qué beneficio obtendrán al ceder su context. La respuesta de Tristan es que el producto debe hacer que el usuario perciba los beneficios del ciclo virtuoso del context antes de que lo abandone.

\begin{figure}[H]
\centering
\includegraphics[width=0.95\textwidth]{context-privacy-tradeoff.png}
\caption{El intercambio entre context y privacidad: beneficios y riesgos de ceder poco, una cantidad moderada o demasiado.}
\end{figure}

\begin{knowledgebox}{Lectura de la figura: más privacidad no siempre es mejor, y más context tampoco}
Ceder poco context implica poco riesgo, pero el usuario no percibe el valor del producto; ceder una cantidad moderada permite mejores emparejamientos y conexiones; ceder demasiado puede hacer que la situación se salga de control. Un buen producto debe permitir que el usuario vea los beneficios y que controle el alcance de la autorización.
\end{knowledgebox}

\subsection{Almacenamiento local y confianza}

En la entrevista se menciona que los usuarios probablemente siempre confiarán más en que su información privada esencial se guarde localmente. Esto significa que los productos sociales de IA deben diseñar memoria local, autorizaciones revocables, registros transparentes, niveles de privacidad y un comportamiento del doble digital que pueda explicarse. De lo contrario, cuanto más esencial sea el context, menos se atreverá el usuario a cederlo.

\subsection{Conexiones reales frente al bullicio de la IA}

Tristan afirma con claridad que tiene poco sentido que las IA publiquen y comenten entre sí. Lo que le importa al usuario es si detrás hay una persona real. La IA puede abrir camino y reducir la barrera para romper el hielo, pero la estrella polar debe ser la tasa de conexiones entre personas reales y la proporción de acciones realizadas por personas reales.

\begin{importantbox}{Métrica estrella polar: tasa de conexiones entre personas reales}
Una red social de IA no existe para que las IA se entretengan entre ellas, sino para que las personas reales se conecten mejor. Por eso, la estrella polar de Elys no es el DAU ni el número de mensajes de IA, sino la tasa de conexiones entre personas reales, la proporción de acciones realizadas por personas reales y si los usuarios encuentran a personas que de verdad valen la pena.
\end{importantbox}

\subsection{Resumen de la sección}

Las redes sociales de IA deben resolver el intercambio entre context y privacidad. Sin context, el doble digital no sirve; sin confianza, el usuario no cede su context; sin conexiones entre personas reales, las redes sociales de IA pierden su sentido.
